%\documentclass{beamer}
%\usepacakge[UTF8]{ctex}
%----------------------------------------------------------------------------------------
%    PACKAGES AND THEMES
%----------------------------------------------------------------------------------------

\documentclass[aspectratio=169,xcolor=dvipsnames]{beamer}
%\usepackage{euler, amsfonts, amssymb}
\usepackage[UTF8]{ctex}
\usetheme{SimplePlus}
\setbeamertemplate{footline}{}

\usepackage{hyperref}
\usepackage{graphicx} % Allows including images
\usepackage{booktabs} % Allows the use of \toprule, \midrule and \bottomrule in tables

\usepackage{mathpartir}

\usepackage{listings}
\usepackage{unicode-math}
\usepackage{fontspec}
%\setmainfont{FreeSerif}
%\setmonofont{FreeMono}
%\setmonofont{Latin Modern Math}
\setmonofont{Unifont}
%\setmonofont{DejaVu Sans Mono}

\newcommand\bbD{\mathbb{D}}
\newcommand\JEq{\mathtt{JEq}}
\newcommand\FTy{\mathtt{FTy}}
\newcommand\comp{\mathtt{comp}}
\newcommand\fst{\mathtt{fst}}
\newcommand\snd{\mathtt{snd}}
\newcommand\Ty{\mathtt{Ty}}
\newcommand\Ter{\mathtt{Ter}}
\newcommand\dM{\mathsf{dM}}
\newcommand\ttp{\mathtt{p}}
\newcommand\ttq{\mathtt{q}}
\newcommand\bfone{\mathbf{1}}
\newcommand\Set{\mathbf{Set}}
\newcommand\DM{\mathbf{DM}}
\newcommand\Id{\mathtt{Id}}
\newcommand\refl{\mathtt{refl}}
\newcommand\C{\mathbf{C}}
\newcommand\psC{\widehat{\mathbf{C}}}
\newcommand\extop{{\scalebox{0.4}[1]{\_}}.{\scalebox{0.8}[1]{\_}}}
\newcommand{\shortminus}{\raisebox{0.3ex}{\scalebox{0.8}[1.0]{-}}}
\newcommand{\mathshortminus}{\mathbin{\text{\normalfont\shortminus}}}
\newcommand{\ul}{{\scalebox{0.4}[1.0]{\_}}}

\DeclareMathOperator\FreeDeMorgan{FreeDeMorgan}
%\DeclareMathOperator\hom{hom}

%----------------------------------------------------------------------------------------
%    TITLE PAGE
%----------------------------------------------------------------------------------------

\title{立方集}
%\subtitle{Subtitle}

%\author{Pin-Yen Huang}

%\institute
%{
%    Department of Computer Science and Information Engineering \\
%    National Taiwan University % Your institution for the title page
%}
%\date{\today} % Date, can be changed to a custom date
\date{}

%----------------------------------------------------------------------------------------
%    PRESENTATION SLIDES
%----------------------------------------------------------------------------------------

\begin{document}

\begin{frame}
    % Print the title page as the first slide
    \titlepage
\end{frame}

\begin{frame}[fragile,shrink]{德摩根代数$(X,\top,\bot,\land,\lor,\neg)$}
\begin{verbatim}
  refl         : ∀ {x} → x ≈ x
  sym          : ∀ {x y} → x ≈ y → y ≈ x
  trans        : ∀ {x y z} → x ≈ y → y ≈ z → x ≈ z
  ∧-cong       : ∀ {x x′ y y′} → x ≈ x′ → y ≈ y′ → (x ∧ y) ≈ (x′ ∧ y′)
  ∨-cong       : ∀ {x x′ y y′} → x ≈ x′ → y ≈ y′ → (x ∨ y) ≈ (x′ ∨ y′)
  ¬-cong       : ∀ {x y} → x ≈ y → (¬ x) ≈ (¬ y)
  ∧-assoc      : ∀ {x y z} → ((x ∧ y) ∧ z) ≈ (x ∧ (y ∧ z))
  ∨-assoc      : ∀ {x y z} → ((x ∨ y) ∨ z) ≈ (x ∨ (y ∨ z)) 
  ∧-comm       : ∀ {x y} → (x ∧ y) ≈ (y ∧ x)
  ∨-comm       : ∀ {x y} → (x ∨ y) ≈ (y ∨ x)
  ∧-idem       : ∀ {x} → (x ∧ x) ≈ x
  ∨-idem       : ∀ {x} → (x ∨ x) ≈ x
  absorb-∧∨    : ∀ {x y} → (x ∧ (x ∨ y)) ≈ x
  absorb-∨∧    : ∀ {x y} → (x ∨ (x ∧ y)) ≈ x
  de-morgan₁   : ∀ {t s} → ¬ (t ∨ s) ≈ (¬ t ∧ ¬ s)
  de-morgan₂   : ∀ {t s} → ¬ (t ∧ s) ≈ (¬ t ∨ ¬ s)
  ¬-involution : ∀ {t} → ¬ (¬ t) ≈ t
  ∨-distrib-∧ˡ : ∀ {t s r} → (t ∨ (s ∧ r)) ≈ ((t ∨ s) ∧ (t ∨ r))
  ∨-distrib-∧ʳ : ∀ {t s r} → (s ∧ r) ∨ t ≈ (s ∨ t) ∧ (r ∨ t)
  ∧-distrib-∨ˡ : ∀ {t s r} → t ∧ (s ∨ r) ≈ ((t ∧ s) ∨ (t ∧ r))
  ∧-distrib-∨ʳ : ∀ {t s r} → (s ∨ r) ∧ t ≈ (s ∧ t) ∨ (r ∧ t)
  ⊤-∧-identity : ∀ {t} → (t ∧ ⊤) ≈ t
  ⊥-∨-identity : ∀ {t} → (t ∨ ⊥) ≈ t
\end{verbatim}
\end{frame}

\begin{frame}[fragile,shrink]{自由德摩根代数}
  \begin{verbatim}
data FreeDeMorgan (X : Set ℓ) : Set ℓ where
  η     : X → FreeDeMorgan X                               
  ⊤     : FreeDeMorgan X                                   
  ⊥     : FreeDeMorgan X                                   
  _∨_   : FreeDeMorgan X → FreeDeMorgan X → FreeDeMorgan X 
  _∧_   : FreeDeMorgan X → FreeDeMorgan X → FreeDeMorgan X 
  ¬_    : FreeDeMorgan X → FreeDeMorgan X                  
  \end{verbatim}
在 $\FreeDeMorgan X$ 上定义上一页展示的等价关系$\approx$
\pause
\begin{itemize}
\item $\neg\ul$, $\top$, $\bot$ 也分别可以记为 $1\mathshortminus\ul$, $1$, $0$
\item $\DM$ 表示德摩根代数组成的范畴
\item 集合 $I$ 上的自由德摩根代数记为 $\dM(I)$
\end{itemize}
\end{frame}

\begin{frame}{立方范畴$\square$}
  \begin{itemize}
    \item 固定一个无限集合$\bbD$, 我们称它的元素为\emph{名字}, 一般记为$i,j,k,\ldots$
    \item $\square$的对象是$\bbD$的有限子集, 一般记为$I,J,K,\ldots$
    \item 态射$\square(I,J)$是所有的函数$J\to\dM(I)$. 这样的函数和德摩根代数同态$\dM(J)\to\dM(I)$之间有一个双射. $\square$中的态射$f\in\square(I,J)$和$g\in\square(J,K)$的复合是$g:K\to\dM(J)$和$f$对应的同态$\dM(J)\to\dM(I)$在$\Set$中的复合
  \end{itemize}
\end{frame}

\begin{frame}{立方集}
  \begin{itemize}
    \item 我们称立方范畴上的预层为立方集: $\square^{op} \to \Set$
    \item 从立方集范畴$\Set^{\square^{op}}$出发, 根据前面介绍的“从预层范畴到族范畴”, 我们可以得到一个 CwF, 它允许我们解释立方类型论的所有特性
  \end{itemize}
\end{frame}

\begin{frame}{记法}
  \begin{itemize}
    \item 对$I\in\square$, $i\notin I$, 我们把 $I\cup \{i\}$ 简记为 $I,i$
    \item $(i0)$代表态射$I\to I,i$ (它是映射 $I,i\to\dM(I)$), 它把$i$映为$0$, 其余恒等输出
    \item 类似, $(i1)$代表态射$I\to I,i$, 它把$i$映为$1$, 其余恒等输出
    \item $s_i:I,i\to I$ 代表包含 $I\subseteq\dM(I,i)$
  \end{itemize}
\end{frame}
\end{document}